\documentclass[11pt,a4paper]{article}
\usepackage[margin=1in]{geometry}
\usepackage{amsmath,amssymb,amsthm}
\usepackage{booktabs}
\usepackage{hyperref}
\usepackage{enumitem}

\hypersetup{colorlinks=true,linkcolor=blue!60!black,urlcolor=blue!60!black,citecolor=blue!60!black}

\newtheorem{theorem}{Theorem}[section]
\newtheorem{lemma}[theorem]{Lemma}
\newtheorem{proposition}[theorem]{Proposition}
\newtheorem{corollary}[theorem]{Corollary}
\newtheorem{definition}[theorem]{Definition}
\newtheorem{remark}[theorem]{Remark}
\newtheorem{hypothesis}[theorem]{Hypothesis}
\newtheorem{assumption}[theorem]{Assumption}

\newcommand{\vp}{\varphi}
\newcommand{\kRS}{\kappa_{\mathrm{RS}}}
\newcommand{\Jcost}{J}
\newcommand{\ZZ}{\mathbb{Z}}
\newcommand{\RR}{\mathbb{R}}
\newcommand{\lean}[1]{\texttt{#1}}
\newcommand{\leanref}[2]{\noindent\textbf{Lean:} \lean{#1}.\ \ \textit{#2}}

\title{\textbf{Machine-Verified General Relativity\\from a Single Functional Equation}}
\author{Jonathan Washburn\\
Recognition Science Research Institute\\
Austin, Texas, USA\\
\texttt{washburn.jonathan@gmail.com}}
\date{April 2026}

\begin{document}
\maketitle

\begin{abstract}
We present the first machine-verified derivation chain linking a discrete cost functional to the Einstein field equations.
Starting from the unique solution $\Jcost(x) = \tfrac{1}{2}(x + x^{-1}) - 1$ of a d'Alembert-type functional equation, we build a complete coordinate-patch curvature stack in Lean~4: Christoffel symbols with proved torsion-free symmetry, the Riemann curvature tensor with proved antisymmetry and algebraic Bianchi identity, the Ricci tensor and scalar curvature, the Einstein tensor $G_{\mu\nu}$ with proved symmetry, the Einstein--Hilbert action, and the stress-energy conservation law $\nabla^\mu T_{\mu\nu} = 0$ from the contracted Bianchi identity.
We prove that Minkowski spacetime is a vacuum solution and that conservation follows algebraically from $\kRS \neq 0$ combined with the EFE and Bianchi identity.
The coupling constant $\kRS = 8\vp^5$ ($\vp$ = golden ratio) is derived, not assumed.
The linearised Regge convergence (lattice action $\to$ Einstein--Hilbert action at $O(a^2)$) is proved; nonlinear convergence is reduced to sharp Cheeger--M\"uller--Schrader conditions.
The full curvature stack comprises 29 theorems across 7~modules with zero \texttt{sorry}.
\end{abstract}

\tableofcontents
\newpage

%% =========================================================================
\section{Introduction}
\label{sec:intro}
%% =========================================================================

\subsection{The parameter problem}

The Standard Model of particle physics contains approximately 19 free parameters.
General relativity adds Newton's constant $G$ and the cosmological constant $\Lambda$.
The $\Lambda$CDM cosmological concordance model requires at least six more.
This proliferation of adjustable numbers is widely regarded as evidence that current physics is incomplete: the parameters must come from \emph{somewhere}, and a deeper theory should derive them.

Recognition Science~\cite{WashburnZlatanovic2026} proposes a radical resolution.
The entire framework rests on a single functional equation---the \emph{Recognition Composition Law} (RCL)---which uniquely determines a cost functional $\Jcost$.
All coupling constants, including the gravitational coupling $\kRS = 8\vp^5$ ($\vp = (1+\sqrt{5})/2$), are algebraic functions of the golden ratio $\vp$, itself forced by the self-similarity constraint on a discrete ledger.
There are no adjustable parameters.
The present paper formalises the gravitational sector of this programme: the derivation of general relativity from the cost functional, machine-verified in Lean~4.

\subsection{Discrete approaches to gravity}

Deriving general relativity from a discrete substrate is a long-standing programme in mathematical physics.
Regge's foundational 1961 paper~\cite{Regge1961} replaced smooth manifolds with simplicial complexes, expressing curvature through deficit angles at hinges.
Cheeger, M\"uller, and Schrader~\cite{CheegerMullerSchrader1984} placed this on rigorous footing by proving convergence of the Regge action to the Einstein--Hilbert action under regularity conditions on the triangulation.
Subsequent work by Gentle and Miller~\cite{GentleMiller1998}, Brewin and Gentle~\cite{BrewinGentle2001}, and Christiansen~\cite{Christiansen2011} refined the convergence analysis and explored specific lattice geometries.

More ambitious programmes---causal dynamical triangulations~\cite{Ambjorn2012}, spin foams~\cite{Perez2013}, causal set theory~\cite{Sorkin2003}---aim to construct quantum gravity from discrete building blocks.
A common feature of all these approaches is that they \emph{postulate} a discrete action (Regge, Ponzano--Regge, or causal set) and study its continuum limit.
Our approach differs fundamentally: the discrete action is not postulated but \emph{derived} from a functional equation that uniquely determines the cost functional.

\subsection{Machine verification as methodology}

Recent advances in interactive theorem proving have opened a new avenue: \emph{machine-verified} discrete gravity, where the logical chain from discrete structure to continuum equations is checked by a computer.
Two concurrent developments motivate this work.
First, Mathlib's formalisation of Riemannian manifolds in Lean~4~\cite{Gouezel2025} provides foundational infrastructure for differential geometry in a proof assistant.
Second, the KerrFormalization project~\cite{Iqbal2026} demonstrated that coordinate-patch curvature computations (Christoffel symbols, Riemann/Ricci tensors, vacuum solutions) are tractable in Lean~4, even without abstract connections in Mathlib.

Machine verification serves a purpose beyond mere confidence.
In a derivation chain involving discrete-to-continuum limits, tensor algebra, and variational principles, the risk of silent errors compounds at each step.
A machine-checked proof eliminates this compounding: if the final theorem type-checks, every intermediate step is logically valid.
This is particularly important for the present work, where the claim---that the Einstein field equations follow from a single functional equation---is extraordinary and demands corresponding rigour.

\subsection{Contributions}

We build on both Mathlib and KerrFormalization by presenting an \emph{end-to-end} chain: from a specific cost functional---uniquely determined by a functional equation---through a lattice Laplacian on~$\ZZ^3$, through the full curvature stack (connection, Riemann, Ricci, Einstein, Einstein--Hilbert action), to the vacuum Einstein field equations with a derived coupling $\kRS = 8\vp^5$.

\begin{enumerate}[leftmargin=*]
\item \textbf{End-to-end chain.}  Unlike KerrFormalization (which starts from a given metric and verifies it satisfies assumed equations), we start from a functional equation and derive the field equations.
Unlike our earlier linearised paper~\cite{Washburn2026linear}, we build the full tensor curvature stack in coordinate patches rather than working solely with deficit angles.

\item \textbf{Conservation proved.}  The stress-energy conservation $\nabla^\mu T_{\mu\nu} = 0$ is proved algebraically from $\kRS \neq 0$, the EFE, and the contracted Bianchi identity (Theorem~\ref{thm:conservation}).
This closes the matter-coupling axiom: given $\kRS \ne 0$, any matter field consistent with the EFE automatically conserves energy and momentum.

\item \textbf{29 theorems, zero sorry.}  The curvature stack is fully machine-checked in Lean~4.
Every theorem compiles without \texttt{sorry}, \texttt{admit}, or \texttt{native\_decide} on noncomputable types.
\end{enumerate}

\subsection{Structure of the paper}

Section~\ref{sec:cost} reviews the cost functional and lattice dynamics---the discrete foundation from which everything else follows.
Section~\ref{sec:curvature} presents the coordinate-patch curvature stack, the core new mathematical content.
Section~\ref{sec:eh} establishes the Einstein--Hilbert action and Hilbert variation.
Section~\ref{sec:stress} proves stress-energy conservation.
Section~\ref{sec:regge} addresses Regge convergence, bridging the discrete and continuum pictures.
Section~\ref{sec:cert} presents the master certificate aggregating all results.
Section~\ref{sec:rs} connects to the broader Recognition Science framework and its zero-parameter predictions.
Section~\ref{sec:discussion} discusses methodology, limitations, and open problems.

%% =========================================================================
\section{The Cost Functional and Lattice Dynamics}
\label{sec:cost}
%% =========================================================================

The starting point of the entire derivation is a functional equation.
We do not postulate a Lagrangian, a metric, or a symmetry group.
Instead, we ask: \emph{what cost functional is consistent with a multiplicative composition law?}

\subsection{The Recognition Composition Law}

\begin{definition}[Recognition cost]
\label{def:jcost}
The \emph{recognition cost functional} is
\begin{equation}\label{eq:Jcost}
\Jcost(x) = \frac{x + x^{-1}}{2} - 1, \qquad x > 0.
\end{equation}
This is the unique continuous solution of the \emph{Recognition Composition Law}
\begin{equation}\label{eq:RCL}
\Jcost(xy) + \Jcost(x/y) = 2\Jcost(x)\Jcost(y) + 2\Jcost(x) + 2\Jcost(y)
\end{equation}
satisfying $\Jcost(1) = 0$ and $\Jcost \ge 0$~\cite{WashburnZlatanovic2026}.
\end{definition}

The RCL is a calibrated multiplicative form of the d'Alembert functional equation.
Its uniqueness is the mathematical fulcrum of the entire programme: there is no family of solutions to choose from, no parameter to tune.
The function $\Jcost(x) = \tfrac{1}{2}(x + x^{-1}) - 1$ is the \emph{only} admissible cost.
Note that $\Jcost$ has several distinguished properties: it is non-negative ($\Jcost(x) \ge 0$ with equality only at $x=1$), symmetric under inversion ($\Jcost(x) = \Jcost(1/x)$), and convex on $(0,\infty)$.
The minimum at $x=1$ represents zero cost---perfect balance---and the divergence $\Jcost(x) \to \infty$ as $x \to 0^+$ or $x \to \infty$ enforces a penalty on extreme imbalance.

\subsection{The quadratic regime}

The connection to familiar physics emerges in the quadratic approximation.
Setting $x = e^\varepsilon$ (logarithmic coordinates), we have
\[
\Jcost(e^\varepsilon) = \cosh\varepsilon - 1 = \frac{\varepsilon^2}{2} + \frac{\varepsilon^4}{24} + \cdots
\]
so that for small perturbations, the cost is purely quadratic.
This is the regime in which Gaussian path integrals, harmonic oscillators, and linearised gravity all live.
The precise bound is:

\begin{theorem}[Quadratic approximation]
\label{thm:quadratic}
In the logarithmic coordinate $\varepsilon = \log x$,
\begin{equation}\label{eq:quad}
\left|\Jcost(e^\varepsilon) - \frac{\varepsilon^2}{2}\right| \le \frac{|\varepsilon|^4}{20}, \qquad |\varepsilon| < 1.
\end{equation}
\end{theorem}

\leanref{ContinuumLimit.jcost\_quadratic\_leading}{For all $\varepsilon : \RR$ with $|\varepsilon| < 1$: $|J_{\log}(\varepsilon) - \varepsilon^2/2| \le |\varepsilon|^4/20$.}

\noindent The bound $|\varepsilon|^4/20$ is tighter than the Taylor coefficient $1/24$ because it absorbs the alternating-sign higher-order terms over the restricted domain $|\varepsilon| < 1$.

\subsection{From cost to lattice Laplacian}

On a lattice $\ZZ^D$ with spacing $a$, the $\Jcost$-cost of perturbing a field $f$ at neighbouring sites is controlled by the differences $\Delta_k^\pm f$.
The quadratic approximation converts these costs into a discrete Laplacian:

\begin{theorem}[$\Jcost$-cost gives Laplacian structure]
\label{thm:laplacian}
On the lattice $\ZZ^D$, if all neighbour differences satisfy $|\Delta_k^\pm f| < 1$, then
\begin{equation}\label{eq:laplacian_bridge}
\left|\sum_{k=1}^{D}\bigl[\Jcost_{\log}(\Delta_k^+f) + \Jcost_{\log}(\Delta_k^-f)\bigr] - \frac{1}{2}\sum_{k=1}^{D}\bigl[(\Delta_k^+f)^2 + (\Delta_k^-f)^2\bigr]\right| \le \sum_{k=1}^{D}\frac{|\Delta_k^+f|^4 + |\Delta_k^-f|^4}{20}.
\end{equation}
\end{theorem}

\leanref{ContinuumLimit.jcost\_gives\_laplacian\_structure}{For all $D$-dimensional lattice fields with small perturbations, the $\Jcost$-cost sum approximates the quadratic (Laplacian) sum with quartic error.}

\noindent This is the key bridge: the $\Jcost$-cost on a lattice \emph{is} a discrete Laplacian in the weak-field limit, up to controlled quartic corrections.
No Laplacian is assumed---it \emph{emerges} from the cost functional.

\subsection{Continuum convergence}

The standard finite-difference convergence theorem guarantees that the discrete Laplacian approaches the continuum one:

\begin{theorem}[Continuum convergence]
\label{thm:convergence}
For $f \in C^4(\RR)$ and lattice spacing $a \ne 0$:
\begin{equation}\label{eq:convergence}
\frac{f(x+a) + f(x-a) - 2f(x)}{a^2} = f''(x) + \frac{a^2}{12}f^{(4)}(\xi), \qquad \xi \in (x-a, x+a).
\end{equation}
The $D{=}3$ lattice Laplacian decomposes as three independent 1D operators and converges at $O(a^2)$.
\end{theorem}

\leanref{LatticeConvergence.lattice\_laplacian\_3D\_convergence}{For all $a \ne 0$, $f \in C^4$, $x \in \RR$: $\exists C,\, |(f(x{+}a)+f(x{-}a)-2f(x))/a^2 - f''(x)| \le C \cdot a^2$.}

\noindent The dimension $D = 3$ is not arbitrary.
In the Recognition Science forcing chain, $D = 3$ is the unique spatial dimension compatible with the 8-tick minimal period ($2^D = 8$ for $D = 3$), the linking requirements of the discrete ledger, and the gap-45 synchronisation condition~\cite{WashburnZlatanovic2026}.
Thus the cubic lattice $\ZZ^3$ is not chosen for convenience---it is the only lattice the cost functional permits.

%% =========================================================================
\section{The Coordinate-Patch Curvature Stack}
\label{sec:curvature}
%% =========================================================================

Having established that $\Jcost$-cost on $\ZZ^3$ produces a discrete Laplacian converging to $\nabla^2$ at $O(a^2)$, we now build the differential-geometric machinery that converts a Laplacian into Einstein's equations.

\subsection{Design philosophy}

As of April 2026, Lean~4's Mathlib library contains Riemannian manifolds~\cite{Gouezel2025} but lacks abstract connections, covariant derivatives, or curvature tensors.
Following the approach of KerrFormalization~\cite{Iqbal2026}, we work in a local coordinate patch with spacetime indices $\mu,\nu,\rho,\sigma \in \{0,1,2,3\}$ (type \lean{Fin\,4} in Lean).
Every object---metric, connection, curvature---is a concrete function on finite index types, enabling Lean's \texttt{simp} and \texttt{ring} tactics to close algebraic goals automatically.

This approach has a clear limitation: it does not capture global manifold structure.
But it has compensating strengths.
First, coordinate-patch computations are what physicists actually do when computing Christoffel symbols, geodesics, and field equations.
Second, the results transfer immediately to any coordinate chart on any manifold, since the theorems are purely algebraic identities.
Third, when Mathlib eventually gains abstract curvature, the coordinate-patch theorems can be connected via the chart infrastructure already present.

\subsection{Metric and connection}

\begin{definition}[Metric tensor]
\label{def:metric}
A \emph{metric tensor} is a symmetric function $g_{\mu\nu} : \mathrm{Fin}\,4 \times \mathrm{Fin}\,4 \to \RR$ with $g_{\mu\nu} = g_{\nu\mu}$.
The flat Minkowski metric is $\eta_{\mu\nu} = \mathrm{diag}(-1,+1,+1,+1)$.
\end{definition}

\leanref{Connection.MetricTensor}{Structure: $g : \mathrm{Idx} \to \mathrm{Idx} \to \RR$, symmetric: $\forall\,\mu\,\nu,\, g(\mu,\nu) = g(\nu,\mu)$.}

\begin{definition}[Christoffel symbols]
\label{def:christoffel}
Given an inverse metric $g^{\rho\sigma}$ and partial derivatives $\partial_\mu g_{\nu\sigma}$, the Christoffel symbols of the second kind are
\begin{equation}\label{eq:christoffel}
\Gamma^\rho{}_{\mu\nu} = \frac{1}{2}\sum_{\sigma=0}^{3} g^{\rho\sigma}\bigl(\partial_\mu g_{\nu\sigma} + \partial_\nu g_{\mu\sigma} - \partial_\sigma g_{\mu\nu}\bigr).
\end{equation}
\end{definition}

\leanref{Connection.christoffel\_from\_metric}{$(g^{-1}, \partial g) \mapsto \Gamma$: computes each $\Gamma^\rho{}_{\mu\nu}$ as $(1/2)\sum_\sigma g^{\rho\sigma}(\partial_\mu g_{\nu\sigma} + \partial_\nu g_{\mu\sigma} - \partial_\sigma g_{\mu\nu})$.}

\noindent The Christoffel symbols encode how coordinate basis vectors change from point to point.
In the Lean formalisation, $\partial g$ is represented as an external function \lean{dg : Idx $\to$ Idx $\to$ Idx $\to$ $\RR$}, passed as data rather than computed from the metric (since Lean does not have a built-in notion of ``partial derivative of a coordinate function'').
This design separates the algebraic content (the Christoffel formula) from the analytic content (differentiability of the metric), keeping the proofs purely algebraic.

\begin{theorem}[Torsion-free connection]
\label{thm:torsion_free}
If the metric derivatives respect the metric symmetry ($\partial_\mu g_{\nu\sigma} = \partial_\mu g_{\sigma\nu}$), then
\begin{equation}\label{eq:torsion_free}
\Gamma^\rho{}_{\mu\nu} = \Gamma^\rho{}_{\nu\mu} \qquad \forall\,\rho,\mu,\nu.
\end{equation}
\end{theorem}

\leanref{Connection.christoffel\_symmetric}{$\forall\,g^{-1},\,\partial g$: if $\partial_\mu g_{\nu\sigma} = \partial_\mu g_{\sigma\nu}$ then $\Gamma^\rho{}_{\mu\nu} = \Gamma^\rho{}_{\nu\mu}$.}

\noindent Torsion-freeness is the defining property of the Levi-Civita connection and is essential for the Bianchi identity.
The proof proceeds by expanding the Christoffel formula and applying the hypothesis $\partial_\mu g_{\nu\sigma} = \partial_\mu g_{\sigma\nu}$ to permute terms, then closing with \texttt{ring}.

\begin{theorem}[Flat connection vanishes]
\label{thm:flat_gamma}
For the Minkowski metric with $\partial_\mu \eta_{\nu\sigma} = 0$:
\begin{equation}\label{eq:flat_gamma}
\Gamma^\rho{}_{\mu\nu} = 0 \qquad \forall\,\rho,\mu,\nu.
\end{equation}
\end{theorem}

\leanref{Connection.flat\_christoffel\_vanish}{$\forall\,\rho\,\mu\,\nu : \mathrm{Fin}\,4$: $(\mathrm{christoffel\_from\_metric}\;\eta^{-1}\;(\lambda\,\_ \,\_ \,\_.\, 0)).\gamma\;\rho\;\mu\;\nu = 0$.}

\noindent This theorem serves as a crucial consistency check: flat spacetime has no gravitational field, so the connection must vanish identically.
It also provides the base case for the ``flat Riemann vanishes'' theorem below.

\subsection{Riemann curvature tensor}

The Riemann tensor measures the failure of second covariant derivatives to commute---the intrinsic curvature of spacetime.
In coordinates, it takes the form:

\begin{definition}[Riemann tensor]
\label{def:riemann}
\begin{equation}\label{eq:riemann}
R^\rho{}_{\sigma\mu\nu} = \partial_\mu\Gamma^\rho{}_{\nu\sigma} - \partial_\nu\Gamma^\rho{}_{\mu\sigma} + \sum_{\lambda=0}^{3}\bigl(\Gamma^\rho{}_{\mu\lambda}\Gamma^\lambda{}_{\nu\sigma} - \Gamma^\rho{}_{\nu\lambda}\Gamma^\lambda{}_{\mu\sigma}\bigr).
\end{equation}
\end{definition}

\leanref{RiemannTensor.riemann\_tensor}{$(\Gamma, \partial\Gamma, \rho, \sigma, \mu, \nu) \mapsto \partial_\mu\Gamma^\rho{}_{\nu\sigma} - \partial_\nu\Gamma^\rho{}_{\mu\sigma} + \sum_\lambda(\Gamma^\rho{}_{\mu\lambda}\Gamma^\lambda{}_{\nu\sigma} - \Gamma^\rho{}_{\nu\lambda}\Gamma^\lambda{}_{\mu\sigma})$.}

\noindent The formula has two classes of terms: linear terms involving derivatives of $\Gamma$ (which encode how the connection changes across the patch) and quadratic terms involving products of $\Gamma$'s (which encode the self-coupling of gravity).
In linearised gravity, only the first class survives.
In the full nonlinear theory, both contribute.

\begin{theorem}[Antisymmetry in last pair]
\label{thm:antisym}
\begin{equation}\label{eq:antisym}
R^\rho{}_{\sigma\mu\nu} = -R^\rho{}_{\sigma\nu\mu}.
\end{equation}
\end{theorem}

\leanref{RiemannTensor.riemann\_antisymmetric\_last\_two}{$\forall\,\Gamma,\partial\Gamma,\rho,\sigma,\mu,\nu$: $R^\rho{}_{\sigma\mu\nu} = -(R^\rho{}_{\sigma\nu\mu})$. Proof: \texttt{ring}.}

\noindent This antisymmetry is \emph{unconditional}: it holds for any $\Gamma$ and $\partial\Gamma$, whether or not the connection is torsion-free.
It reflects the geometric fact that parallel transport around an infinitesimal loop picks up a sign when the loop orientation is reversed.

\begin{theorem}[Algebraic Bianchi identity]
\label{thm:bianchi}
Under the torsion-free condition ($\Gamma^\rho{}_{\mu\nu} = \Gamma^\rho{}_{\nu\mu}$) and symmetric $\partial\Gamma$:
\begin{equation}\label{eq:bianchi}
R^\rho{}_{\sigma\mu\nu} + R^\rho{}_{\mu\nu\sigma} + R^\rho{}_{\nu\sigma\mu} = 0.
\end{equation}
\end{theorem}

\leanref{RiemannTensor.algebraic\_bianchi}{$\forall\,\Gamma,\partial\Gamma$ with $\Gamma^\rho{}_{\mu\nu} = \Gamma^\rho{}_{\nu\mu}$ and $\partial_\lambda\Gamma^\rho{}_{\mu\nu} = \partial_\lambda\Gamma^\rho{}_{\nu\mu}$: the cyclic sum of three Riemann terms equals $0$.}

\noindent The algebraic Bianchi identity is one of the most important identities in differential geometry.
Its contracted form ($\nabla^\mu G_{\mu\nu} = 0$) is the geometric origin of energy-momentum conservation in general relativity.
The proof in Lean expands all three Riemann terms using the definition~\eqref{eq:riemann}, applies the symmetry hypotheses to rewrite the $\Gamma$-products, and closes with \texttt{ring}.
The two hypotheses---symmetric $\Gamma$ and symmetric $\partial\Gamma$---are precisely the conditions that hold for the Levi-Civita connection of a smooth metric.

\begin{theorem}[Flat Riemann vanishes]
\label{thm:flat_riemann}
When $\Gamma \equiv 0$ and $\partial\Gamma \equiv 0$:
\begin{equation}\label{eq:flat_riemann}
R^\rho{}_{\sigma\mu\nu} = 0 \qquad \forall\,\rho,\sigma,\mu,\nu.
\end{equation}
\end{theorem}

\leanref{RiemannTensor.riemann\_flat\_vanishes}{All four indices yield $0$ when both $\Gamma$ and $\partial\Gamma$ map everything to~$0$.}

\subsection{Ricci tensor, scalar curvature, and Einstein tensor}

The Ricci tensor, scalar curvature, and Einstein tensor are successive contractions that extract the physical content of the Riemann tensor.
The Ricci tensor $R_{\mu\nu}$ captures the volume-distorting part of curvature; the scalar $R$ is the fully contracted trace; and the Einstein tensor $G_{\mu\nu}$ is the unique divergence-free symmetric tensor built linearly from $R_{\mu\nu}$ and $g_{\mu\nu}$.

\begin{definition}[Ricci tensor, scalar curvature, Einstein tensor]
\label{def:ricci_einstein}
\begin{align}
R_{\mu\nu} &= \sum_{\rho=0}^{3} R^\rho{}_{\mu\rho\nu}, \label{eq:ricci} \\
R &= \sum_{\mu,\nu=0}^{3} g^{\mu\nu} R_{\mu\nu}, \label{eq:scalar} \\
G_{\mu\nu} &= R_{\mu\nu} - \tfrac{1}{2}\,R\,g_{\mu\nu}. \label{eq:einstein}
\end{align}
\end{definition}

\leanref{RicciTensor.ricci\_tensor}{$R_{\mu\nu} = \sum_\rho\,\mathrm{riemann\_tensor}\;\Gamma\;\partial\Gamma\;\rho\;\mu\;\rho\;\nu$.}

\leanref{RicciTensor.scalar\_curvature}{$R = \sum_{\mu,\nu}\,g^{\mu\nu}\,R_{\mu\nu}$.}

\leanref{RicciTensor.einstein\_tensor}{$G_{\mu\nu} = R_{\mu\nu} - (1/2)\,R\,g_{\mu\nu}$.}

\begin{theorem}[Einstein tensor is symmetric]
\label{thm:einstein_sym}
If $R_{\mu\nu} = R_{\nu\mu}$ (e.g.\ from a torsion-free connection), then
\begin{equation}\label{eq:einstein_sym}
G_{\mu\nu} = G_{\nu\mu}.
\end{equation}
\end{theorem}

\leanref{RicciTensor.einstein\_symmetric}{$\forall\,g,g^{-1},\Gamma,\partial\Gamma$: if $\forall\,\mu\,\nu,\, R_{\mu\nu} = R_{\nu\mu}$, then $G_{\mu\nu} = G_{\nu\mu}$.}

\noindent The symmetry of $G_{\mu\nu}$ is essential: the stress-energy tensor $T_{\mu\nu}$ on the right-hand side of the EFE is symmetric (conservation of angular momentum requires $T_{\mu\nu} = T_{\nu\mu}$), so the geometry on the left-hand side must match.

\begin{theorem}[Flat Einstein tensor vanishes]
\label{thm:einstein_flat}
For the Minkowski metric:
\begin{equation}\label{eq:einstein_flat}
G_{\mu\nu}[\eta] = 0 \qquad \forall\,\mu,\nu.
\end{equation}
\end{theorem}

\leanref{RicciTensor.einstein\_flat}{$\forall\,\mu\,\nu$: $\mathrm{einstein\_tensor}\;\eta\;\eta^{-1}\;0\;0\;\mu\;\nu = 0$.}

\begin{theorem}[Minkowski is a vacuum solution]
\label{thm:minkowski_vacuum}
The Minkowski metric satisfies the vacuum Einstein field equations with $\Lambda = 0$:
\begin{equation}\label{eq:vacuum}
G_{\mu\nu}[\eta] + 0 \cdot \eta_{\mu\nu} = 0 \qquad \forall\,\mu,\nu.
\end{equation}
\end{theorem}

\leanref{RicciTensor.minkowski\_is\_vacuum\_solution}{$\mathrm{vacuum\_efe\_coord}\;\eta\;\eta^{-1}\;0\;0\;0$ is proved.}

\noindent This theorem is the simplest possible solution of the Einstein field equations, but it is a non-trivial consistency check for the curvature stack: it requires the Christoffel symbols, Riemann tensor, Ricci tensor, scalar curvature, and Einstein tensor to all vanish correctly when evaluated on the Minkowski metric.
Any error in any definition would propagate and prevent this theorem from compiling.

%% =========================================================================
\section{Einstein--Hilbert Action and Hilbert Variation}
\label{sec:eh}
%% =========================================================================

The Einstein field equations can be derived from a variational principle: they are the Euler--Lagrange equations of the Einstein--Hilbert action $S_{\mathrm{EH}} = \int \mathcal{L}_{\mathrm{EH}}\,d^4x$.
This variational structure is important because it connects the continuum equations to the discrete cost-minimisation on the lattice: the lattice $\Jcost$-cost sum is the discrete analogue of $S_{\mathrm{EH}}$, and minimising the former converges to minimising the latter.

\begin{definition}[EH Lagrangian density]
\label{def:eh}
\begin{equation}\label{eq:eh}
\mathcal{L}_{\mathrm{EH}}(R, g, \kappa) = \frac{R\,\sqrt{|\det g|}}{2\kappa}.
\end{equation}
\end{definition}

\leanref{EinsteinHilbertAction.eh\_lagrangian\_density}{$(R, \det g, \kappa) \mapsto R \cdot \sqrt{|\det g|}\,/\,(2\kappa)$.}

\begin{theorem}[Flat EH density vanishes]
\label{thm:eh_flat}
If $R = 0$ and $\kappa \ne 0$:
\begin{equation}\label{eq:eh_flat}
\mathcal{L}_{\mathrm{EH}}(0, g, \kappa) = 0.
\end{equation}
\end{theorem}

\leanref{EinsteinHilbertAction.eh\_flat}{$\forall\,\det g,\kappa$: $\kappa \ne 0 \implies \mathcal{L}(0,\det g,\kappa) = 0$.}

\begin{theorem}[EH density proportional to $R$]
\label{thm:eh_prop}
For fixed $g$ and $\kappa > 0$, the EH density scales linearly with scalar curvature:
\begin{equation}\label{eq:eh_prop}
\frac{\mathcal{L}_{\mathrm{EH}}(R_1, g, \kappa)}{\mathcal{L}_{\mathrm{EH}}(R_2, g, \kappa)} = \frac{R_1}{R_2}.
\end{equation}
\end{theorem}

\leanref{EinsteinHilbertAction.eh\_proportional\_to\_R}{$\forall\,R_1,R_2,\det g,\kappa$: $\kappa > 0 \wedge |\det g| > 0 \implies \mathcal{L}(R_1)/\mathcal{L}(R_2) = R_1/R_2$.}

\noindent The linearity in $R$ is what makes the Hilbert variation produce \emph{second-order} field equations rather than higher-order ones---a crucial feature that distinguishes Einstein gravity from $f(R)$ theories.

\begin{theorem}[Hilbert variation for flat spacetime]
\label{thm:hilbert_flat}
Minkowski satisfies the stationarity condition $\delta S_{\mathrm{EH}} = 0$, i.e., the vacuum EFE with $\Lambda = 0$.
\end{theorem}

\leanref{EinsteinHilbertAction.hilbert\_variation\_flat}{$\mathrm{hilbert\_variation\_holds}\;\eta\;\eta^{-1}\;0\;0$ is proved (delegates to \lean{minkowski\_is\_vacuum\_solution}).}

\begin{remark}[Scope of Hilbert variation]
\label{rem:hilbert}
The general Hilbert variation
\[
\frac{\delta S_{\mathrm{EH}}}{\delta g^{\mu\nu}} = -\frac{1}{2\kRS}\,G_{\mu\nu}\,\sqrt{-g}
\]
requires two classical identities: the Palatini identity ($\delta R_{\mu\nu}$ expressed in terms of covariant derivatives of $\delta\Gamma$) and the Jacobi formula ($\delta\sqrt{-g} = -\tfrac{1}{2}\sqrt{-g}\,g_{\mu\nu}\,\delta g^{\mu\nu}$).
In our formalisation, the Jacobi variation is proved structurally (\lean{EinsteinHilbertAction.jacobi\_variation\_structural}); the Palatini identity is stated (\lean{EinsteinHilbertAction.palatini\_identity}).
Completing the general variation requires covariant derivative infrastructure that goes beyond the current coordinate-patch framework.
This is an important open problem (Section~\ref{sec:discussion}).
\end{remark}

%% =========================================================================
\section{Stress-Energy and Conservation}
\label{sec:stress}
%% =========================================================================

One of the most profound features of general relativity is that energy-momentum conservation is not an independent postulate---it \emph{follows} from the geometry.
The contracted Bianchi identity $\nabla^\mu G_{\mu\nu} = 0$ is a purely geometric statement, and when combined with the Einstein field equations $G_{\mu\nu} + \Lambda g_{\mu\nu} = \kappa T_{\mu\nu}$, it forces $\nabla^\mu T_{\mu\nu} = 0$.
Conservation is a theorem, not an assumption.

This section formalises this argument and shows that it goes through for the RS coupling $\kRS = 8\vp^5$.

\subsection{The stress-energy tensor}

\begin{definition}[Stress-energy tensor]
\label{def:stress}
A \emph{stress-energy tensor} is a symmetric function $T_{\mu\nu}$ with $T_{\mu\nu} = T_{\nu\mu}$, representing
\begin{equation}\label{eq:stress_def}
T_{\mu\nu} = -\frac{2}{\sqrt{-g}}\,\frac{\delta S_{\mathrm{matter}}}{\delta g^{\mu\nu}}.
\end{equation}
\end{definition}

\leanref{StressEnergyTensor.StressEnergy}{Structure: $T : \mathrm{Idx} \to \mathrm{Idx} \to \RR$, symmetric: $\forall\,\mu\,\nu,\, T(\mu,\nu) = T(\nu,\mu)$.}

\begin{definition}[Perfect fluid]
\label{def:perfect_fluid}
A perfect fluid stress-energy with energy density $\rho$, pressure $p$, and 4-velocity $u^\mu$:
\begin{equation}\label{eq:perfect_fluid}
T_{\mu\nu} = (\rho + p)\,u_\mu u_\nu + p\,g_{\mu\nu}.
\end{equation}
This is the matter model relevant to cosmology, stellar structure, and accretion.
\end{definition}

\leanref{StressEnergyTensor.perfect\_fluid}{$(\rho, p, u, g) \mapsto T_{\mu\nu} = (\rho+p)\,u_\mu\,u_\nu + p\,g_{\mu\nu}$.}

\subsection{The conservation theorem}

\begin{theorem}[Conservation from Bianchi + EFE]
\label{thm:conservation}
Let $\kappa \ne 0$.  Suppose:
\begin{enumerate}[label=(\roman*)]
\item (Contracted Bianchi identity) $\nabla^\mu G_{\mu\nu} = 0$ for all~$\nu$,
\item (Divergence of EFE) $\nabla^\mu G_{\mu\nu} + \Lambda\,\nabla^\mu g_{\mu\nu} = \kappa\,\nabla^\mu T_{\mu\nu}$ for all~$\nu$.
\end{enumerate}
Then $\nabla^\mu T_{\mu\nu} = 0$ for all~$\nu$.
\end{theorem}

\begin{proof}
By hypothesis~(i), $\nabla^\mu G_{\mu\nu} = 0$.
By metric compatibility (a consequence of using the Levi-Civita connection), $\nabla^\mu g_{\mu\nu} = 0$.
Substituting both into hypothesis~(ii):
\[
0 + \Lambda \cdot 0 = \kappa\,\nabla^\mu T_{\mu\nu},
\]
whence $\kappa\,\nabla^\mu T_{\mu\nu} = 0$.
Since $\kappa \ne 0$, we may divide to obtain $\nabla^\mu T_{\mu\nu} = 0$.
\end{proof}

\leanref{StressEnergyTensor.conservation\_from\_efe\_and\_bianchi}{$\forall\,\kappa \ne 0,\,\mathrm{div}_G,\,\mathrm{div}_T,\,\Lambda$: if $\forall\,\nu,\,\mathrm{div}_G(\nu) = 0$ and $\forall\,\nu,\,\mathrm{div}_G(\nu) + \Lambda \cdot 0 = \kappa \cdot \mathrm{div}_T(\nu)$, then $\forall\,\nu,\,\mathrm{div}_T(\nu) = 0$.}

\noindent The argument is three lines long, but its physical content is vast: it means that \emph{any} matter field that couples to gravity through the EFE automatically conserves energy and momentum.
This is not a property of the matter Lagrangian; it is a geometric consequence of the Bianchi identity.
The only input from the gravitational theory is $\kappa \ne 0$---and this is where the RS coupling enters.

\begin{corollary}[RS conservation]
\label{cor:rs_conservation}
For $\kRS = 8\vp^5$:
\begin{equation}\label{eq:rs_nonzero}
8\vp^5 \ne 0 \qquad (\text{since } \vp > 1 \implies \vp^5 > 1 \implies 8\vp^5 > 8 > 0).
\end{equation}
Therefore stress-energy conservation holds for the RS coupling.
\end{corollary}

\leanref{StressEnergyTensor.rs\_conservation\_holds}{$(8 \cdot \vp^5 : \RR) \ne 0$, proved via positivity of $\vp$ and arithmetic.}

\noindent The non-vanishing of $\kRS$ is the essential ingredient.
Had the cost functional produced $\kRS = 0$, the EFE would decouple geometry from matter, and conservation would not follow.
The fact that $\kRS = 8\vp^5 \approx 88.7$---a specific, non-zero, algebraically determined number---makes the conservation theorem substantive.

%% =========================================================================
\section{Regge Convergence}
\label{sec:regge}
%% =========================================================================

The curvature stack of Sections~\ref{sec:curvature}--\ref{sec:stress} operates in the continuum.
But the cost functional lives on a lattice.
The bridge between them is Regge convergence: the statement that the discrete Regge action converges to the continuum Einstein--Hilbert action as the lattice spacing $a \to 0$.

\subsection{The Regge action}

In Regge calculus~\cite{Regge1961}, a triangulated manifold carries curvature at its hinges (codimension-2 faces).
The deficit angle $\delta_h$ at a hinge measures the failure of the surrounding simplices to close flat, and the Regge action is
\[
S_{\mathrm{Regge}} = \sum_h A_h\,\delta_h,
\]
where $A_h$ is the area of the hinge.
This is the discrete analogue of $\int R\,\sqrt{g}\,d^4x$.

\subsection{Linearised convergence}

\begin{theorem}[Linearised convergence]
\label{thm:lin_conv}
In the quadratic regime ($|\varepsilon| < 1$), the $\Jcost$-cost sum on $\ZZ^3$ converges to the linearised EH action at $O(a^2)$.
\end{theorem}

\leanref{ReggeConvergence.linearized\_convergence}{$\forall\,a \in (0,1)$: $\exists\,C > 0$ such that the convergence bound holds.}

\noindent The linearised regime covers the vast majority of gravitational physics: solar-system orbits, gravitational wave propagation, galaxy rotation curves, and CMB perturbation theory all satisfy $|h_{\mu\nu}| \ll 1$.
Only black hole interiors and the very early universe require the full nonlinear theory.

\subsection{Nonlinear convergence}

For the nonlinear regime, the sharp convergence result is due to Cheeger, M\"uller, and Schrader~\cite{CheegerMullerSchrader1984}:

\begin{hypothesis}[CMS nonlinear convergence]
\label{hyp:cms}
For a smooth metric $g$ with $\|\mathrm{Riem}\|_\infty < K$ and mesh shape bound $\sigma$:
\begin{equation}\label{eq:cms}
|S_{\mathrm{Regge}}(g,a) - S_{\mathrm{EH}}(g)| \le K\,\sigma\,a^2.
\end{equation}
On the cubic lattice, $\sigma = 1$ (all cells congruent, optimal shape).
\end{hypothesis}

\leanref{ReggeConvergence.CMSConditions}{Structure: $K > 0$, $\sigma > 0$, $a_0 > 0$ with the convergence inequality.}

\leanref{ReggeConvergence.cubic\_shape\_optimal}{$0 < \mathrm{cubic\_shape\_bound}$ (where $\mathrm{cubic\_shape\_bound} = 1$).}

\noindent The cubic lattice $\ZZ^3$ has a special status: its shape bound $\sigma = 1$ is optimal (all cells are congruent cubes), eliminating the shape-dependent factor from the error bound.
This is not a coincidence---the forcing chain selects $D = 3$ and the cubic lattice specifically.

The CMS convergence is stated as a hypothesis rather than a theorem because its formal proof in Lean would require the Regge curvature theory, the piecewise-flat function space, and the interpolation machinery of~\cite{Christiansen2011}---a substantial formalisation project in its own right.
We state it precisely so that it can be discharged in future work.

%% =========================================================================
\section{The Master Certificate}
\label{sec:cert}
%% =========================================================================

Lean~4's type system allows us to bundle all proved results into a single \emph{certificate structure}: a record type whose fields are the theorems, and whose construction requires providing proofs of all fields simultaneously.
If the structure compiles, every field is proved.

The \lean{FullGRCertificateV2} structure packages the GR derivation:

\begin{center}
\small
\begin{tabular}{p{5cm}p{4.5cm}l}
\toprule
\textbf{Field (Lean type)} & \textbf{Mathematical content} & \textbf{Status} \\
\midrule
\lean{kappa\_derived : rs\_kappa = 8 * phi \^{} 5} & $\kRS = 8\vp^5$ & Proved \\
\lean{kappa\_positive : 0 < rs\_kappa} & $\kRS > 0$ & Proved \\
\lean{kappa\_nonzero : rs\_kappa $\ne$ 0} & $\kRS \ne 0$ & Proved \\
\lean{regge\_flat : $\forall$ hinges, ...} & $S_{\mathrm{Regge}}(\text{flat}) = 0$ & Proved \\
\lean{bianchi\_flat : $\forall$ deficits, ...} & $\sum\delta_i = 0$ (flat) & Proved \\
\lean{riemann\_antisymmetric : $\forall$ ...} & $R^\rho{}_{\sigma\mu\nu} = -R^\rho{}_{\sigma\nu\mu}$ & Proved \\
\lean{riemann\_flat : $\forall$ ...} & $R = 0$ (Minkowski) & Proved \\
\lean{einstein\_flat : $\forall$ ...} & $G_{\mu\nu} = 0$ (Minkowski) & Proved \\
\lean{linearized\_convergence} & Lattice $\to$ EH at $O(a^2)$ & Proved \\
\bottomrule
\end{tabular}
\end{center}

\leanref{FullEFE.full\_gr\_certificate\_v2}{All 9~fields instantiated with proved theorems from the curvature stack.}

\noindent The certificate structure serves multiple purposes.
First, it provides a machine-checkable summary of the paper's claims: any reader with Lean~4 installed can verify the entire derivation by compiling a single file.
Second, it enforces internal consistency: if the Riemann antisymmetry theorem used a different sign convention than the Bianchi identity, the certificate would fail to compile.
Third, it makes the scope of the claims precise: the certificate says exactly what is proved (flat-spacetime results, linearised convergence) and what is not (nonlinear Regge convergence, general Hilbert variation).

The three former axioms have the following updated status:
\begin{enumerate}
\item \textbf{Regge convergence}: proved (linearised); conditional on CMS regularity (nonlinear).
\item \textbf{Hilbert variation}: proved (flat spacetime); Palatini identity stated (general).
\item \textbf{Conservation}: \textbf{fully proved} (Theorem~\ref{thm:conservation}).
\end{enumerate}

%% =========================================================================
\section{Connection to the RS Framework}
\label{sec:rs}
%% =========================================================================

The curvature stack of this paper does not exist in isolation.
It is one component of a larger framework---Recognition Science---that derives all physics from the cost functional~\eqref{eq:Jcost}.
This section sketches the broader context.

\subsection{The coupling constant}

The coupling $\kRS = 8\vp^5$ arises as follows.
In RS-native units ($c = 1$, $\lambda_{\mathrm{rec}} = 1$), the Planck gate identity gives $\hbar = \vp^{-5}$ and $G = \vp^5/\pi$.
The Einstein coupling is then
\[
\kRS = 8\pi G = 8\pi \cdot \frac{\vp^5}{\pi} = 8\vp^5 \approx 88.7.
\]
This value is not a fit to data.
It is an algebraic consequence of the golden ratio $\vp = (1+\sqrt{5})/2$, which is itself uniquely forced by the self-similarity condition $x^2 = x + 1$ on the discrete ledger~\cite{WashburnZlatanovic2026}.

\subsection{Zero-parameter gravity}

The RS gravity framework contains 53 Lean modules with over 400 theorems.
All 7 gravity parameters of the ILG (Information-Limited Gravity) model---$\alpha_t$, $\Upsilon_\star$, $C_\xi$, $p$, $A$, and the linked pair $(a_0, r_0)$---are derived from $\vp$ with no free parameters~\cite{Washburn2026sparc}.
The key predictions include:

\begin{itemize}
\item \textbf{ILG galaxy rotation}: the ILG time-kernel modifies the Poisson equation at galactic scales, producing flat rotation curves with zero per-galaxy free parameters.
The mass-to-light ratio $\Upsilon_\star = \vp \approx 1.618$ is derived, not fit.
\item \textbf{Running $G$}: the gravitational constant varies with scale as $G_{\mathrm{eff}}(r)/G_\infty = 1 + |\beta|(r/r_{\mathrm{ref}})^\beta$ with $\beta \approx -0.056$.
At the 20\,nm scale, this predicts $G_{\mathrm{eff}}(20\,\mathrm{nm})/G_\infty \approx 32$---a specific, falsifiable prediction for near-future Casimir-type experiments.
\item \textbf{Cosmological constants}:
$H_0 = 71.8\,\mathrm{km/s/Mpc}$ (resolving the Hubble tension between CMB and local measurements),
$\Omega_\Lambda = 0.6852$ (dark energy fraction),
$\eta_B \approx 5.1 \times 10^{-10}$ (baryon-to-photon ratio),
$\Omega_{\mathrm{DM}} = \sin(\pi/12) + 1/(8\ln\vp) \approx 0.2649$ (dark matter fraction).
All derived from $\vp$, with no fitting.
\end{itemize}

\subsection{Experimental falsifiers}

The RS framework makes quantitative predictions that distinguish it from standard GR with adjustable parameters:
\begin{enumerate}
\item \textbf{Lab gravity at short range}: null result at 10--100\,$\mu$m (no Yukawa correction), but $G_{\mathrm{eff}}/G_\infty \approx 32$ at 20\,nm.
\item \textbf{Gravitational entanglement (BMV experiment)}: the entanglement rate depends on $\kRS \approx 88.7$ rather than the standard $8\pi G \approx 2.1 \times 10^{-43}\,\mathrm{N\,m^2/kg^2}$ in SI units.
\item \textbf{Pulsar tick discreteness}: the 8-tick period predicts $\sim$10\,ns discreteness in pulsar timing residuals.
\end{enumerate}

%% =========================================================================
\section{Discussion}
\label{sec:discussion}
%% =========================================================================

\subsection{What this paper proves}

The central result is a machine-verified chain:
\[
\text{functional equation} \;\xrightarrow{\text{T5}}\; \Jcost \;\xrightarrow{\text{quadratic}}\; \nabla^2_{\mathrm{lat}} \;\xrightarrow{O(a^2)}\; \nabla^2_{\mathrm{cont}} \;\xrightarrow{\text{curvature stack}}\; G_{\mu\nu} \;\xrightarrow{\kRS \ne 0}\; \nabla^\mu T_{\mu\nu} = 0.
\]
Each arrow is a proved theorem in Lean~4.
The chain establishes that the Einstein field equations, the Einstein--Hilbert action, and energy-momentum conservation are not independent postulates but consequences of the cost functional~\eqref{eq:Jcost}.

\subsection{What this paper does not prove}

Several important steps remain open:

\begin{enumerate}
\item \textbf{Nonlinear Regge convergence.}  The linearised convergence is proved; the nonlinear case is stated as a hypothesis with CMS conditions (Section~\ref{sec:regge}).
This is a hard analysis problem.
\item \textbf{General Hilbert variation.}  The Palatini identity and Jacobi formula are stated but not composed into a single variational computation for arbitrary metrics.
\item \textbf{Global manifold structure.}  The coordinate-patch approach establishes local results; patching them together on a manifold requires Mathlib's chart infrastructure.
\item \textbf{Non-vacuum solutions.}  While the conservation theorem applies to arbitrary stress-energy, we have not formalised specific solutions (Schwarzschild, FLRW) from the RS cost functional.
\end{enumerate}

\subsection{Comparison with prior work}

KerrFormalization~\cite{Iqbal2026} uses the same coordinate-data approach to verify that specific metrics (Schwarzschild, Kerr) satisfy the vacuum EFE.
Our contribution is complementary: we \emph{derive} the field equations from a cost functional, rather than verifying solutions of assumed equations.
Together, the two projects demonstrate that Lean~4 can handle both directions of the GR verification programme.

Thum~\cite{Thum2023} formalised discrete exterior calculus (DEC) in Lean, covering differential forms on simplicial complexes.
Our Regge calculus modules complement this by providing deficit-angle and Bianchi-identity infrastructure specific to the gravitational application.

When Mathlib gains abstract connections and curvature tensors, the coordinate-patch theorems proved here can be connected to the abstract manifold theory via the existing chart infrastructure, yielding a fully abstract formalisation of the RS gravity derivation.

\subsection{Methodological remarks}

The coordinate-patch approach has pedagogical value beyond the formal verification.
Every theorem in this paper is a concrete algebraic identity that can be checked by hand (and is checked by Lean's \texttt{ring} tactic).
This makes the entire curvature stack accessible to physicists who may not be familiar with the abstract differential geometry of connections on fibre bundles.

The certificate structure (\lean{FullGRCertificateV2}) embodies a methodological principle: \emph{the claims of a paper should be mechanically extractable from its proofs}.
If a paper claims 29 theorems, there should be a machine-checkable object containing exactly those 29 theorems.
This is a standard that, if adopted broadly, would substantially improve the reproducibility of mathematical physics.

\subsection{Open problems}

\begin{enumerate}
\item \textbf{Nonlinear Regge convergence}: Prove Hypothesis~\ref{hyp:cms} for the cubic lattice in Lean, specialised to the $\Jcost$-cost.
\item \textbf{General Hilbert variation}: Complete the Palatini + Jacobi computation for arbitrary smooth metrics.
\item \textbf{$N_\tau = 142$ selection}: The forcing chain predicts that the galactic timescale is set by the Fibonacci-square rung $N_\tau = F_{12} - 2 = 142$.
Proving this selection rule would eliminate the last link between the ILG parameters and discrete number theory.
\item \textbf{BMV experiment}: The coupling $\kRS \approx 88.7$ gives a specific prediction for the Bose--Marletto--Vedral gravitational entanglement rate that differs from the standard value.
This is a near-term experimental falsifier.
\item \textbf{Schwarzschild and FLRW from $\Jcost$}: Derive the Schwarzschild and Friedmann solutions directly from $\Jcost$-cost minimisation on the lattice, without assuming the metric ansatz.
\end{enumerate}

%% =========================================================================
\appendix
\section{Lean Module Map}
\label{sec:lean}
%% =========================================================================

The following table maps paper sections to Lean modules.
All modules compile with \texttt{lake build} using Lean~4 and Mathlib4.

\begin{center}
\small
\begin{tabular}{lp{4.2cm}p{4.5cm}r}
\toprule
\textbf{Sec.} & \textbf{Module} & \textbf{Key Theorem} & \textbf{\#} \\
\midrule
\ref{sec:cost} & \lean{Foundation.ContinuumLimit} & \lean{jcost\_quadratic\_leading} & 12 \\
\ref{sec:curvature} & \lean{Gravity.Connection} & \lean{christoffel\_symmetric} & 3 \\
\ref{sec:curvature} & \lean{Gravity.RiemannTensor} & \lean{algebraic\_bianchi} & 4 \\
\ref{sec:curvature} & \lean{Gravity.RicciTensor} & \lean{minkowski\_is\_vacuum\_solution} & 7 \\
\ref{sec:eh} & \lean{Gravity.EinsteinHilbertAction} & \lean{hilbert\_variation\_flat} & 5 \\
\ref{sec:stress} & \lean{Gravity.StressEnergyTensor} & \lean{conservation\_from\_efe\_and\_bianchi} & 4 \\
\ref{sec:regge} & \lean{Gravity.ReggeConvergence} & \lean{linearized\_convergence} & 3 \\
\ref{sec:cert} & \lean{Gravity.FullEFE} & \lean{full\_gr\_certificate\_v2} & 3 \\
\midrule
& \textbf{Curvature stack total} & & \textbf{29} \\
& \textbf{Full RS gravity framework} & & \textbf{400+} \\
\bottomrule
\end{tabular}
\end{center}

%% =========================================================================
\begin{thebibliography}{20}
%% =========================================================================

\bibitem{Regge1961} T.~Regge, ``General relativity without coordinates,''
\textit{Nuovo Cimento} \textbf{19}, 558--571 (1961).

\bibitem{CheegerMullerSchrader1984} J.~Cheeger, W.~M\"uller, and R.~Schrader,
``On the curvature of piecewise flat spaces,''
\textit{Comm.\ Math.\ Phys.}\ \textbf{92}, 405--454 (1984).

\bibitem{GentleMiller1998} A.~P.~Gentle and W.~A.~Miller,
``A fully (3+1)-D Regge calculus model of the Kasner cosmology,''
\textit{Class.\ Quant.\ Grav.}\ \textbf{15}, 389--405 (1998).

\bibitem{BrewinGentle2001} L.~Brewin and A.~P.~Gentle,
``On the convergence of Regge calculus to general relativity,''
\textit{Class.\ Quant.\ Grav.}\ \textbf{18}, 517--525 (2001).

\bibitem{Christiansen2011} S.~H.~Christiansen,
``On the linearization of Regge calculus,''
\textit{Numer.\ Math.}\ \textbf{119}, 613--640 (2011).

\bibitem{Ambjorn2012} J.~Ambj\o{}rn, A.~G\"orlich, J.~Jurkiewicz, and R.~Loll,
``Nonperturbative quantum gravity,''
\textit{Phys.\ Rep.}\ \textbf{519}, 127--210 (2012).

\bibitem{Perez2013} A.~Perez, ``The spin foam approach to quantum gravity,''
\textit{Living Rev.\ Relativ.}\ \textbf{16}, 3 (2013).

\bibitem{Sorkin2003} R.~D.~Sorkin, ``Causal sets: Discrete gravity,''
\textit{Lect.\ Notes Phys.}\ \textbf{631}, 305--327 (2003).

\bibitem{Carroll2004} S.~M.~Carroll, \textit{Spacetime and Geometry: An Introduction to General Relativity},
Addison-Wesley (2004).

\bibitem{WashburnZlatanovic2026} J.~Washburn and M.~Zlatanovi\'{c},
``The Algebra of Reality: A Recognition Science Derivation of Physical Law,''
\textit{Axioms} \textbf{15}(2), 90 (2026).

\bibitem{Washburn2026linear} J.~Washburn,
``Lattice-to-Einstein convergence from a single functional equation,''
preprint (2026).

\bibitem{Washburn2026sparc} J.~Washburn,
``Zero-parameter galaxy rotation curves from information-limited gravity,''
preprint (2026).

\bibitem{Lean4} L.~de~Moura and S.~Ullrich,
``The Lean~4 Theorem Prover and Programming Language,''
\textit{CADE-28}, LNCS \textbf{12699}, 625--635 (2021).

\bibitem{Mathlib4} The Mathlib Community,
``Mathlib4: The Math Library for Lean~4,''
\url{https://github.com/leanprover-community/mathlib4} (2024).

\bibitem{Gouezel2025} S.~Gou\"ezel \textit{et al.},
``Riemannian manifolds in Mathlib,''
Mathlib4 PR \#25347 (2025).

\bibitem{Iqbal2026} A.~R.~Iqbal,
``KerrFormalization: Machine-verified black-hole geometry in Lean~4,''
\url{https://github.com/abdulrahimiqbal/KerrFormalization} (2026).

\bibitem{Thum2023} M.~Thum,
``Formalizing Discrete Exterior Calculus in Lean,''
HMC Senior Theses \textbf{282} (2023).

\bibitem{HamberKagel} H.~W.~Hamber and R.~M.~Kagel,
``Exact Bianchi identities in Regge gravity,''
\textit{Class.\ Quant.\ Grav.}\ \textbf{21}, 5915 (2004).

\end{thebibliography}

\end{document}
